%%% This file is automatically generated by LeanArchitect. %%%

%%% Macro definitions for \inputleannode, \inputleanmodule %%%

\makeatletter

% \newleannode{name}{latex} defines a new Lean node
\providecommand{\newleannode}[2]{%
  \expandafter\gdef\csname leannode@#1\endcsname{#2}}
% \inputleannode{name} inputs a Lean node
\providecommand{\inputleannode}[1]{%
  \csname leannode@#1\endcsname}

% \newleanmodule{module}{latex} defines a new Lean module
\providecommand{\newleanmodule}[2]{%
  \expandafter\gdef\csname leanmodule@#1\endcsname{#2}}
% \inputleanmodule{module} inputs a Lean module
\providecommand{\inputleanmodule}[1]{%
  \csname leanmodule@#1\endcsname}

\makeatother

%%% Start of main content %%%

\newleannode{MeromorphicOnRectangle}{\input{/home/runner/work/trough/trough/library/.lake/build/blueprint/module/PrimeNumberTheoremAnd/GeneralMeromorphic.artifacts/MeromorphicOnRectangle}}

\newleanmodule{PrimeNumberTheoremAnd.GeneralMeromorphic}{
\inputleannode{MeromorphicOnRectangle}

A meromorphic function has a pole of finite order.
\begin{definition}\label{PoleOrder}
If $f$ has a pole at $z_0$, then there is an integer $n$ such that
$$
\lim_{z\to z_0} (z-z_0)^n f(z) = c \neq 0.
$$
\end{definition}


If a meromorphic function $f$ has a pole at $z_0$, then the residue of
$f$ at $z_0$ is the coefficient of $1/(z-z_0)$ in the Laurent series of
$f$ around $z_0$.
\begin{definition}\label{Residue}
If $f$ has a pole of order $n$ at $z_0$, then
$$
Res_{z_0} f =
  \lim_{z\to z_0}\frac1{(n-1)!}
    (\partial/\partial z)^{n-1}[(z-z_0)^{n-1}f(z)].
$$
\end{definition}


We can evaluate a small integral around a pole by taking the residue.


If a function $f$ is meromorphic at $z_0$ with a pole of order $n$, then
the residue at $z_0$ of the logarithmic derivative is $-n$ exactly.

}